\documentclass[12pt,a4paper]{article}
\usepackage[utf8]{inputenc}
\usepackage[T1]{fontenc}
\usepackage[spanish,es-noquoting]{babel}
\usepackage{geometry}
\usepackage{xcolor}
\usepackage{enumitem}

\geometry{a4paper, margin=1in}

\title{Evaluación de Examen Teórico-Práctico}
\author{Prof. César Rodríguez \\ \small Curso: Seguridad Informática}
\date{\today}

\begin{document}

\maketitle

\section*{Datos del Estudiante}
\textbf{Nombre:} Juan Álvarez \\
\textbf{Grupo:} 2 (OSINT y Huella Digital) \\
\textbf{Calificación Final:} \textbf{\textcolor{blue}{9.5/10}}

\section*{Evaluación Detallada}

\begin{itemize}[label=$\bullet$]
    \item \textbf{Pregunta 1 (Significado de OSINT):} \textbf{2/2 puntos}. Define correctamente OSINT (Inteligencia de fuentes abiertas) y aporta un excelente enfoque al mencionar que se hace ``sin penetrar la seguridad'', lo que subraya el carácter no intrusivo (pasivo) de esta técnica.
    
    \item \textbf{Pregunta 2 (Datos críticos en WHOIS):} \textbf{2/2 puntos}. Identifica acertadamente las tres piezas clave de información buscadas en la fase de reconocimiento: fechas (registro y modificación), servidores DNS y el dueño del dominio (contactos).
    
    \item \textbf{Pregunta 3 (Google Dork \texttt{site:dominio.com -www}):} \textbf{2/2 puntos}. Comprende el propósito de la exclusión. Explica que al omitir el sitio principal, se expone el resto de la infraestructura (subdominios) evitando ``ruido'' en los resultados.
    
    \item \textbf{Pregunta 4 (Búsqueda de archivos \texttt{.ini} o \texttt{.env}):} \textbf{1.5/2 puntos}. La lógica de agrupamiento usando paréntesis y el operador \texttt{|} (equivalente lógico de \texttt{OR}) es muy buena. Sin embargo, introduce un espacio inválido en el operador (\texttt{site: ejemplo.com}) que rompe la consulta.
    
    \item \textbf{Pregunta 5 (Error HTTP 429 en Python):} \textbf{2/2 puntos}. Respuesta sobresaliente. Reconoce el error por exceso de peticiones y propone dotar al código de resiliencia: limitar el ritmo de escaneo mediante pausas y atrapar el error para guardar la información recolectada hasta ese punto.
\end{itemize}

\vspace{1cm}

\section*{Comentarios Finales}
\noindent El estudiante demuestra un entendimiento sólido tanto de la teoría de reconocimiento pasivo como de las prácticas de programación defensiva para automatizar estas tareas. Posee un enfoque analítico excelente, aunque debe cuidar los detalles estrictos de sintaxis en el buscador.

\vspace{0.5cm}
\begin{center}
    \textit{¡Excelente trabajo!}
\end{center}

\end{document}